%! Confidence Intervals for the Difference of Two Proportions — Unit 6, Lesson 6.8 — Warm-Up
\documentclass[10pt]{article}
\usepackage{apstats-article}
\usepackage{apstats-boxes}
\newcommand{\TallMath}[1]{$\displaystyle #1\rule[-1.4em]{0pt}{3.2em}$}

\begin{document}

\pageheader{Unit 6, Lesson 6.8}{Warm-Up}
\namedateperiod

\vspace{0.15in}

A sports analyst wants to compare the proportion of athletes who use a recovery supplement
in two leagues. League A: $n_1 = 150$, $\hat p_1 = 0.48$. League B: $n_2 = 120$, $\hat p_2 = 0.35$.

\begin{enumerate}
  \item What is the point estimate of $p_1 - p_2$?
        $\hat p_1 - \hat p_2 = $ \blank{3cm}

  \item Verify the Large Counts condition for each group ($\ge 10$ for all four values):
        \begin{itemize}
          \item $n_1 \hat p_1 = $ \blank{3cm} $\ge 10$? \blank{1cm} \quad $n_1(1-\hat p_1) = $ \blank{3cm} $\ge 10$? \blank{1cm}
          \item $n_2 \hat p_2 = $ \blank{3cm} $\ge 10$? \blank{1cm} \quad $n_2(1-\hat p_2) = $ \blank{3cm} $\ge 10$? \blank{1cm}
        \end{itemize}

  \item In a one-sample $z$-interval for $p$, we have one $SE$ term. For the two-sample case,
        write the formula for $SE_{\hat p_1 - \hat p_2}$:
        \[SE = \sqrt{\dfrac{\blank{3cm}}{n_1} + \dfrac{\blank{3cm}}{n_2}}\]

  \item Compute $SE$ for this example.
        $SE = $ \blank{6cm} $\approx$ \blank{2cm}
\end{enumerate}

\end{document}
